\section{Determinante angular y relaci\'on electrod\'ebil de Weinberg}

El determinante de \(T\) es

\begin{equation}
\det T=q_+q_-=D_A.
\label{eq:detT}
\end{equation}

El sector constitutivo del vac\'io eval\'ua \(F(T)^2\), donde
\(F(z)=(1-z^{90})/(1-z^{120})\), mientras que el sector angular
electrod\'ebil eval\'ua \(\det T\). Se trata de dos funciones de un
antecedente com\'un y no de una factorizaci\'on horizontal entre constantes.
Esta distinci\'on se desarrolla en el \cref{ch:ew}.

\begin{remark}[Control negativo]
El bloque queda refutado si se viola cualquiera de estas identidades exactas:
positividad de \(I-T^{120}\), igualdad
\cref{eq:three-four-completion}, relaciones \cref{eq:vacuum-relations} o
compatibilidad \(E_m(\mathcal V_{m+1})=\mathcal V_m\). Una coincidencia SI
no puede compensar un fallo operatorio.
\end{remark}

\subsection*{Alcance lógico y dominio de validez}
Operador, espectro, completaci\'on \(3\to4\), paridad, amplificaci\'on y secciones constitutivas: \emph{Teorema interno exacto}. La formulaci\'on que los re\'une en un solo operador es \emph{Formalizaci\'on nueva}; los datos angulares y las rectas son \emph{Resultado antecedente demostrado} del desarrollo HMT precedente y se declaran aqu\'i como estructura de partida expl\'icita.

